\chapter{Introduction}

\section{Background}

Randomized Controlled Trials (RCTs) remain the gold standard for evaluating the efficacy and safety of novel therapeutic interventions. 
However, the clinical development sector is increasingly constrained by rising operational costs, strict time limitations, and complex 
ethical considerations \citep{dimasi_innovation_2016}. For instance, recruiting a sufficient number of patients, particularly in case of 
rare diseases, poses a major obstacle to the development of new treatments. Furthermore, patients are sometimes randomized to a placebo arm rather than to an 
active control arm, which raises significant ethical concerns regarding patient welfare when other effective treatments already exist \citep{freedman_scientific_1987, world_medical_association_world_2025}.

Historically, the dominant statistical paradigm in the design of RCTs has been the frequentist approach. This framework evaluates long-run operating 
characteristics, relying heavily on strict Type I Error (T1E rate) control. A fundamental limitation of this paradigm is that it treats every trial 
in isolation, ignoring pre-existing information such as data from previous clinical phases or historical control arms that could otherwise 
help the current study. 

Conversely, Bayesian statistical frameworks offer a mathematically rigorous solution to this challenge by incorporating external evidence 
through informative prior distributions. Integrating background knowledge during the planning phase can lead to highly efficient designs, 
accelerating trial execution, reducing operational waste, and limiting the number of concurrent patients required for the control arm.

This not only decreases the overall sample size but also makes the trial more attractive to potential participants, who face a lower 
probability of being assigned to a placebo group. As noted by O'Hagan et al., \textit{"the design of a trial 
is ultimately the sponsor’s decision and its outcome is the sponsor’s risk; therefore, the design should incorporate all available information to 
maximize the probability of a satisfactory outcome"} \citep{ohagan_assurance_2005}.
Recognizing these practical benefits, regulatory bodies are increasingly supportive of these methods, as 
evidenced by recent guidance on the use of Bayesian methodology in clinical trials \citep{fda_bayesian_2026}.

 

\section{Bayesian Designs}
The Bayesian framework formally incorporates historical information into the design through prior distributions, that can be placed either directly 
on the treatment effect or on one (or both) arms:

\begin{itemize}
        \item \textbf{Borrowing on the Control Group:} when recruiting patients for a standard-of-care or placebo arm is ethically challenging or 
    operationally slow, historical data can be used to inform the control parameters. Designs that supplement a smaller concurrent control group
    with historical information are formally referred to in the literature as \textit{hybrid-control} trials \citep[e.g.,][]{gravestock2017adaptive, gravestock2019power, li2022hybrid, fu2023covariate, kopp2024simulating}.


    \item \textbf{Borrowing on the Treatment Effect:} When the primary goal is to use existing efficacy data to support a new indication or
    demographic, such as leveraging adult data to pediatric cohorts, the prior can be placed on the treatment contrast directly \citep{psioda_bayesian_2019}.
\end{itemize}





Evaluating the operating characteristics of Bayesian designs requires particular care. Traditional frequentist designs rely on
conditional power, the probability of success under a single, fixed treatment effect. The Bayesian framework instead shifts the focus
to unconditional power, or \textit{assurance}, the expected power averaged over a prior distribution for the unknown treatment effect.
Integrating over a prior gives a more realistic assessment of the predictive probability of a
statistically significant result \citep{spiegelhalter1993bayesian, ohagan_assurance_2005}. It is also known as the probability of study success \citep{wang_evaluating_2013}
or Bayesian \textit{predictive power} \citep{spiegelhalter_monitoring_1986}.

This thesis focuses on borrowing historical information on the control arm. In this setting, traditional \textit{assurance} is not the relevant metric, since
it averages over uncertainty in the treatment effect, not over the control parameter. Evaluating such designs instead requires metrics 
that account for uncertainty in the control arm. The classical frequentist T1E rate, while commonly used, cannot be strictly
controlled once external information is borrowed \citep{viele_use_2014, kopp2024simulating}. This is acknowledged in regulatory guidance: when a
frequentist T1E rate is not applicable to a Bayesian design that borrows external information, alternative trial characteristics should be
considered \citep{fda_bayesian_2026}. To address this, \citet{best_beyond_2025}
propose the average T1E rate, which integrates the conditional T1E rate over a design prior on the control parameter rather than
 evaluating it at a single point. This reformulation is attractive because it restores the
guarantee that borrowing removes: when the design prior is aligned with the analysis prior, the average T1E rate is
controlled exactly at the nominal level.


\section{Dynamic Borrowing and Prior-Data Conflict}
\label{subsec:info-borrowing}

While borrowing historical data presents clear advantages, it is not without risk. Historical and current trial data may be heterogeneous
due to shifts in demographics, standard of care, or measurement techniques \citep{evans_checking_2006}. If a model forces the combination
of historical and concurrent data when such a discrepancy exists, it can lead to biased estimates and inflated error rates.

Recognizing these dangers, \citet{pocock_combination_1976} proposed the first formal criteria for combining historical and concurrent
controls in a randomized design: the historical trial should be comparable in patient population, treatment regimen, and endpoint
definition, and conducted under sufficiently similar conditions. Several Bayesian frameworks have since been developed to achieve a more
principled form of discounting, in which historical information is downweighted in proportion to heterogeneity.  The common goal of all
these \textit{dynamic borrowing} methods is to borrow strength when historical data align with the current trial
while discounting when a true discrepancy, commonly called \textit{drift} \citep{viele_use_2014}, is present. Examples include
\textit{power priors} \citep{ibrahim_power_2000}, which discount historical data via a scalar power parameter $a_0 \in [0,1]$, and \textit{commensurate priors} \citep{hobbs_2011}, which link historical and current parameters through a shared commensurability model.

Other than those, the \textit{meta-analytic-predictive} (MAP) approach \citep{neuenschwander2010summarizing} derives the prior for the control arm
from a hierarchical meta-analysis of historical trials, treating the trial-specific control data as exchangeable. The MAP prior is
not available in closed form but can be approximated by a mixture of conjugate distributions using simulation methods. However, standard 
MAP prior can still react too slowly when historical and current data disagree. Therefore, it is possible to augment the mixture with a weakly informative
component whose weight $\omega$ encodes the prior probability that the new trial is not exchangeable with the historical controls.
This \textit{robust MAP} prior responds more rapidly to \textit{prior-data conflict}, shifting posterior weight toward the vague component as
divergence between historical and current data grows. This thesis focuses primarily on the MAP framework and its robust extension,
an approach that is genuinely used in practice. A scoping review of Phase II trials identified 37 studies applying robust
Bayesian borrowing, of which $56.8\%$ borrowed in order to supplement the control arm, and roughly a fifth of those
borrowing from external sources did so through robust mixture priors \citep{weru2026uptake}. Recent examples include
\citet{richeldi2022trial} and \citet{chen2024efficacy}. That same review also found that the parameters governing the
degree of borrowing are poorly reported, with only a third of the trials applying a fixed downweighting explaining how
the weight had been chosen, which motivates a systematic study of how that weight shapes the operating characteristics.



\section{Dual Criterion} 
\label{subsec:dual_criterion_intro}

A limitation of standard trial designs is that success is defined by a single statistical threshold: a significant $p$-value or a posterior probability. While this guarantees sufficient evidence against the null hypothesis, 
it says nothing about whether the effect is large enough to be clinically meaningful, since a statistically significant result can arise from a modest, practically unimportant treatment difference.

The dual criterion responds to this by requiring two conditions rather than one. Developed for proof of concept
trials, where the decision is whether to continue development \citep{gsponer_practical_2014, fisch2015bayesian}, it
answers two distinct questions. The first asks whether the drug works at all, and demands high posterior confidence that the effect
is superior to placebo. The second asks whether the drug has a relevant clinical effect, and demands only moderate
posterior confidence (usually 50\%) that the effect exceeds a target difference, the threshold below which a treatment is not worth developing further. 

The two conditions together admit four outcomes. Both holding
gives a GO and neither gives a NO-GO, while the two mixed cases are indeterminate and call for different
responses. Evidence of a relevant effect that misses the significance threshold often suggests expanding the trial, whereas an
effect that is significant but marginal may suggests reconsidering the target population.


\citet{fisch2015bayesian} further show that an informative prior built from historical placebo data can substitute for a
large part of the trial, with a design of twelve treated and six control patients performing comparably to one of forty
patients split evenly. 


\citet{roychoudhury_beyond_2018} later formalise this as the \textit{dual criterion} (DC) design, in which trial success requires both a statistical criterion and a clinical criterion: the estimated effect must reach a pre-specified Decision Value (DV), 
the minimum clinically meaningful improvement elicited from domain experts. power at the DV is approximately $50\%$ 
by construction, meaning the sample size must be chosen to achieve adequate power for effects meaningfully above the DV.




\section{Sequential Designs}
\label{subsec:sequential_intro}

A classical fixed design decides in advance for a single analysis at a final predetermined sample size. Interim analyses relax this constraint and allow the trial to stop earlier. The motivation is both ethical and financial. Stopping early for efficacy makes a working
treatment available sooner and saves patients from an inferior arm, while stopping early for futility avoids exposing further patients to an ineffective one and saves resources for other research.

Historically, interim analyses have been applied predominantly within the frequentist framework,
where taking multiple looks at the data inflates the T1E rate beyond its nominal level. Because the frequentist test statistic depends only 
on the difference between the two arms, its behaviour across looks can be represented as a discretely observed Brownian 
motion \citep{jennison_group_1999}. Various methods to define boundaries to control the T1E rate have been developed, including Pocock's constant threshold 
\citep{pocock_group_1977}, O'Brien and Fleming's boundaries \citep{obrien_multiple_1979}, and the alpha-spending function \citep{lan_demets_discrete_1983}.

Bayesian group sequential designs replace the frequentist test statistic with the posterior probability that the treatment is superior, 
stopping early for efficacy once this probability is high enough, and for futility using posterior or predictive probability. Traditionally, 
frequentist operating characteristics are obtained using simulation approaches \citep{berry2010bayesian}.
Whether these thresholds should be adjusted for the planning of interim analyses to guarantee frequentist T1E control is itself debated, 
and we refer the reader to \citet{zhou_ji_2024} for a complete overview. \citet{ryan2020we} give an empirical answer: efficacy 
stopping inflates the T1E rate as more looks are added unless the thresholds are adjusted, while futility stopping alone does not, 
since it can only remove trials that would have failed regardless.

\citet{zhu2017bayesian} propose an algorithm to calibrate the Bayesian thresholds to a target T1E rate, and \citet{shi2019control} 
generalize this method to other design types and add a theoretical shortcut that avoids simulation under a simple conjugate prior. 

An interim analysis can also be used to assess \textit{prior-data conflict}. If external information does not agree with the current trial,
more patients may be randomized to the control arm after the interim, or the randomization ratio may be adjusted \citep{schmidli_robust_2014}.


\section{The Experimental Unit Information Index}
\label{subsec:euii_intro}
Comparing designs is not straightforward when they differ in sample size, power, and T1E rate simultaneously, especially when
the design is adaptive, i.e. the sample size is random. \citet{held2025balancing} address this with the Experimental Unit
Information Index (EUII), which normalizes the \textit{diagnostic odds ratio} (DOR) by the number of patients a design needs, making
designs with different stopping rules comparable on a common, per-patient scale.


\section{Objectives \note{TO DO}}
\note{This section is built up progressively as the thesis is revised. Add an objective here once the corresponding chapter is settled.}

The objectives of this thesis are the following.

\begin{itemize}
    \item Extend the Bayesian fixed design metrics to group sequential designs with robust mixture priors. This covers the
    average T1E rate of \citet{best_beyond_2025} together with an analogous average power. The latter replaces
    \textit{assurance}, which averages over uncertainty in the treatment effect rather than over the control parameter and
    is therefore not the relevant quantity when historical information is borrowed on the control arm.

    \item Use the EUII to evaluate whether the sequential framework delivers higher evidentiary value per unit than the
    fixed design framework.
\end{itemize}


Look at the discussion of the paper: Beyond Randomized Clinical Trials: Use of External Controls
Quite nice

The remainder of the thesis is organized as follows. Chapter~\ref{sec:methodology} develops the methodology: the MAP prior, the average metrics with their closed forms under Normal priors, the extension of the EUII to the borrowing framework, the \textit{dual criterion}, and group sequential designs. Chapter~\ref{sec:Applications} applies the framework to a case study in Crohn's disease, evaluating fixed and sequential designs under aligned and misspecified design priors. Chapter~\ref{sec:conclusion} concludes with a discussion of the findings and directions for future work.



%%% Local Variables: 
%%% mode: latex
%%% TeX-master: "../MasterThesisSfS"
%%% End: 
